\chapter{环形计数器变式题（15题完整解析）}

% ============================================================
\section{题1：4位环形计数器}

\textbf{题目}：设计4位环形计数器。初始1000。状态：1000→0100→0010→0001→1000。

\begin{solution}[完整解答]

\textbf{【结构】}4位移位寄存器，Q0反馈回D3（右移+反馈）。

\textbf{【逐拍分析】}
\begin{center}
\begin{tabular}{|c|c|c|c|c|}\hline
拍 & Q3 & Q2 & Q1 & Q0 \\ \hline
0(初) & 1 & 0 & 0 & 0 \\ \hline
1 & 0 & 1 & 0 & 0 \\ \hline
2 & 0 & 0 & 1 & 0 \\ \hline
3 & 0 & 0 & 0 & 1 \\ \hline
4 & 1 & 0 & 0 & 0（循环！）\\ \hline
\end{tabular}
\end{center}

\textbf{【VHDL】}
\begin{lstlisting}
entity ring4 is
  port (clk, rst_n : in std_logic; q : out std_logic_vector(3 downto 0));
end ring4;

architecture behavioral of ring4 is
  signal reg : std_logic_vector(3 downto 0);
begin
  process(clk, rst_n)
  begin
    if rst_n = '0' then reg <= "1000";  -- One-hot初始
    elsif rising_edge(clk) then
      reg <= reg(0) & reg(3 downto 1);  -- 右移+反馈
    end if;
  end process;
  q <= reg;
end behavioral;
\end{lstlisting}

\textbf{【为什么是4个状态】}4个D触发器，只有1个1在循环移动。1在4个位置→4个状态。

\textbf{【关键VHDL语句】}\texttt{reg <= reg(0) \& reg(3 downto 1)}
\begin{itemize}
  \item reg(0)→reg(3)：最低位反馈到最高位（闭环！）
  \item reg(3)→reg(2)：原最高位右移一位
  \item reg(2)→reg(1)：右移
  \item reg(1)→reg(0)：右移
\end{itemize}

\end{solution}


% ============================================================
\section{题2：8位环形计数器}

\textbf{题目}：设计8位环形计数器（8个状态）。

\begin{solution}[完整解答]

\textbf{【VHDL】}
\begin{lstlisting}
entity ring8 is
  port (clk, rst_n : in std_logic; q : out std_logic_vector(7 downto 0));
end ring8;

architecture behavioral of ring8 is
  signal reg : std_logic_vector(7 downto 0);
begin
  process(clk, rst_n)
  begin
    if rst_n = '0' then reg <= "10000000";
    elsif rising_edge(clk) then
      reg <= reg(0) & reg(7 downto 1);  -- 8位环形
    end if;
  end process;
  q <= reg;
end behavioral;
\end{lstlisting}

\textbf{【状态循环】}$10000000\to01000000\to\cdots\to00000001\to10000000$

共8个状态（8拍一个循环）。

\end{solution}


% ============================================================
\section{题3：4位约翰逊计数器（扭环形）}

\textbf{题目}：设计4位约翰逊计数器。初始0000。状态序列：0000→1000→1100→1110→1111→0111→0011→0001→0000。

\begin{solution}[完整解答]

\textbf{【反馈】}$\overline{Q_0}$→D3（反相最低位反馈回最高位）。

\textbf{【逐拍分析】}
\begin{center}
\begin{tabular}{|c|c|c|c|c|}\hline
拍 & Q3 & Q2 & Q1 & Q0 \\ \hline
0 & 0 & 0 & 0 & 0 \\ \hline
1 & 1 & 0 & 0 & 0 \\ \hline
2 & 1 & 1 & 0 & 0 \\ \hline
3 & 1 & 1 & 1 & 0 \\ \hline
4 & 1 & 1 & 1 & 1 \\ \hline
5 & 0 & 1 & 1 & 1 \\ \hline
6 & 0 & 0 & 1 & 1 \\ \hline
7 & 0 & 0 & 0 & 1 \\ \hline
8 & 0 & 0 & 0 & 0（循环！）\\ \hline
\end{tabular}
\end{center}

\textbf{【VHDL】}
\begin{lstlisting}
entity johnson4 is
  port (clk, rst_n : in std_logic; q : out std_logic_vector(3 downto 0));
end johnson4;

architecture behavioral of johnson4 is
  signal reg : std_logic_vector(3 downto 0);
begin
  process(clk, rst_n)
  begin
    if rst_n = '0' then reg <= (others => '0');
    elsif rising_edge(clk) then
      reg <= (not reg(0)) & reg(3 downto 1);  -- 反相反馈!
    end if;
  end process;
  q <= reg;
end behavioral;
\end{lstlisting}

\textbf{【关键VHDL】}\texttt{reg <= (not reg(0)) \& reg(3 downto 1)}

注意 \texttt{not reg(0)} — 反相！这是约翰逊与普通环形计数器的唯一区别。

\textbf{【状态数】}4个DFF→$2\times4=8$个状态。比普通环形计数器多一倍。

\end{solution}


% ============================================================
\section{题4：环形计数器 vs 二进制计数器}

\textbf{题目}：对比4状态环形计数器与2位二进制计数器。

\begin{solution}[完整解答]

\begin{center}
\begin{tabular}{|l|c|c|}\hline
\textbf{特性} & \textbf{4位环形(4状态)} & \textbf{2位二进制(4状态)} \\ \hline
DFF数量 & 4 & 2 \\
状态编码 & One-hot(1000,0100,...) & Binary(00,01,10,11) \\
解码门数 & 0（读Qx即可）& 需要2-4译码器 \\
关键路径延迟 & 1级DFF($t_{cko}$) & $t_{cko}$+进位链+$t_{setup}$ \\
最大频率 & 最高(D触发器本身限制) & 较低(进位链限制) \\
面积 & 4×DFF & 2×DFF + 加法器 \\
\hline
\end{tabular}
\end{center}

\textbf{【选择规则】}需要高速时选环形计数器；面积受限时选二进制计数器。

\end{solution}


% ============================================================
\section{题5：自启动环形计数器}

\textbf{题目}：4位环形计数器可能进入非法状态。设计自启动逻辑。

\begin{solution}[完整解答]

\textbf{【非法状态检测】}合法状态：one-hot（恰好一个1）。非法状态：全0、多个1。

\textbf{【VHDL——自启动】}
\begin{lstlisting}
entity ring4_selfstart is
  port (clk, rst_n : in std_logic; q : out std_logic_vector(3 downto 0));
end ring4_selfstart;

architecture behavioral of ring4_selfstart is
  signal reg : std_logic_vector(3 downto 0);
  function count_ones(v : std_logic_vector) return integer is
    variable c : integer := 0;
  begin
    for i in v'range loop if v(i)='1' then c:=c+1; end if; end loop;
    return c;
  end function;
begin
  process(clk, rst_n)
  begin
    if rst_n = '0' then reg <= "1000";
    elsif rising_edge(clk) then
      if count_ones(reg) /= 1 then  -- 非法状态检测
        reg <= "1000";              -- 强制恢复
      else
        reg <= reg(0) & reg(3 downto 1);
      end if;
    end if;
  end process;
  q <= reg;
end behavioral;
\end{lstlisting}

\end{solution}


% ============================================================
\section{题6：自启动约翰逊计数器}

\textbf{题目}：约翰逊计数器也可能进入非法状态。设计自启动逻辑。

\begin{solution}[完整解答]

\textbf{【分析】}4位约翰逊计数器合法状态=8个。非法状态=8个（全16个状态的一半）。检测不在合法状态中→强制恢复。

\textbf{【VHDL——预存合法状态ROM】}
\begin{lstlisting}
type rom_type is array (0 to 7) of std_logic_vector(3 downto 0);
constant LEGAL : rom_type := (
  "0000","1000","1100","1110","1111","0111","0011","0001"
);
signal is_legal : boolean;
begin
  is_legal <= false;
  for i in 0 to 7 loop
    if reg = LEGAL(i) then is_legal <= true; end if;
  end loop;

  if not is_legal then
    reg <= "0000";  -- 强制恢复
  else
    reg <= (not reg(0)) & reg(3 downto 1);
  end if;
\end{lstlisting}

\end{solution}


% ============================================================
\section{题7：用环形计数器做序列发生器}

\textbf{题目}：用4位环形计数器输出Q3产生序列。Q3序列是什么？

\begin{solution}[完整解答]

\textbf{【逐拍分析Q3】}
\begin{center}
\begin{tabular}{|c|c|}\hline
拍 & Q3 \\ \hline
0 & 1 \\ \hline 1 & 0 \\ \hline 2 & 0 \\ \hline 3 & 0 \\ \hline
4 & 1 \\ \hline 5 & 0 \\ \hline ... & ... \\ \hline
\end{tabular}
\end{center}

Q3输出序列：\textbf{1, 0, 0, 0, 1, 0, 0, 0, ...} — 周期4。

\textbf{【各Q输出对比】}
\begin{itemize}
  \item Q3: 1,0,0,0,1,0,0,0,...（相位0°）
  \item Q2: 0,1,0,0,0,1,0,0,...（相位90°）
  \item Q1: 0,0,1,0,0,0,1,0,...（相位180°）
  \item Q0: 0,0,0,1,0,0,0,1,...（相位270°）
\end{itemize}

四个Q输出正好是相位差90°的4路信号——可用于多相时钟生成。

\end{solution}


% ============================================================
\section{题8：左移环形计数器}

\textbf{题目}：设计左移环形计数器（1从右向左移动）。

\begin{solution}[完整解答]

\textbf{【VHDL】}
\begin{lstlisting}
reg <= reg(2 downto 0) & reg(3);  -- 左移循环
\end{lstlisting}

\textbf{【状态序列(初始0001)】}
$0001\to0010\to0100\to1000\to0001$

1从右向左循环移动。

\end{solution}


% ============================================================
\section{题9：N位环形计数器状态数推导}

\textbf{题目}：N位移位寄存器构建环形计数器，能产生多少个状态？

\begin{solution}[完整解答]

\textbf{【普通环形计数器】}
One-hot编码。初始化后恰好一个1。N个位置→N个状态。

\textbf{答案：N个状态。}

\textbf{【约翰逊计数器】}
反相反馈。N个DFF→2N个不同状态。

\textbf{答案：2N个状态。}

\textbf{【证明（约翰逊）】}
每个DFF在循环中都经历0和1的完整变化。反相反馈使得1连续填入，0连续推出。每N拍状态翻转一次，共2N拍回到起点。

\end{solution}


% ============================================================
\section{题10：5位环形计数器时序控制器}

\textbf{题目}：用5位环形计数器实现5步时序控制器。Q4-Q0分别作为步骤1-5的使能信号。

\begin{solution}[完整解答]

\textbf{【VHDL】}
\begin{lstlisting}
entity seq_controller is
  port (clk, rst_n : in std_logic;
        step : out std_logic_vector(4 downto 0));
end seq_controller;

architecture behavioral of seq_controller is
  signal reg : std_logic_vector(4 downto 0);
begin
  process(clk, rst_n)
  begin
    if rst_n = '0' then reg <= "10000";  -- 步骤1激活
    elsif rising_edge(clk) then
      reg <= reg(0) & reg(4 downto 1);
    end if;
  end process;
  step <= reg;
end behavioral;
\end{lstlisting}

\textbf{【应用】}多周期操作的控制器（如乘法器的状态控制）。每拍恰好一个步骤激活。

\end{solution}


% ============================================================
\section{题11：给定波形判断计数器类型}

\textbf{题目}：4位输出波形：1000→1100→1110→1111→0111→0011→0001→0000→1000→...。判断类型和状态数。

\begin{solution}[完整解答]

\textbf{【分析】}
\begin{itemize}
  \item 状态数=8（不同模式有8种）
  \item 观察规律：1从左向右填充，然后0从右向左清空
  \item 这是典型的约翰逊计数器！
\end{itemize}

\textbf{【结论】}4位约翰逊计数器（8个状态）。

\textbf{【判断方法】}
\begin{itemize}
  \item One-hot（只有一个1循环）= 普通环形
  \item 111...→011...→001...模式 = 约翰逊（扭环形）
  \item 连续递增的二进制值 = 二进制计数器
\end{itemize}

\end{solution}


% ============================================================
\section{题12：LFSR vs 环形计数器}

\textbf{题目}：比较LFSR和环形计数器的异同。

\begin{solution}[完整解答]

\begin{center}
\begin{tabular}{|l|c|c|}\hline
 & LFSR & 环形计数器 \\ \hline
结构 & 移位寄存器+XOR反馈 & 移位寄存器+直接反馈 \\
反馈逻辑 & 多输入异或（抽头）& 单输入（reg(0)或~reg(0)）\\
状态数(N个DFF) & 最多$2^N-1$ & 最多2N \\
状态编码 & 伪随机 & One-hot \\
应用 & 随机数、CRC & 时序控制、多相信号 \\
\hline
\end{tabular}
\end{center}

\end{solution}


% ============================================================
\section{题13：约翰逊计数器做分频器}

\textbf{题目}：4位约翰逊计数器（8状态）。Q3输出频率？

\begin{solution}[完整解答]

\textbf{【分析】}Q3在8拍中：0→1→1→1→1→0→0→0→0→...。周期=8拍。

$$f_{Q3} = f_{clk} / 8$$

\textbf{【各Q频率（4位约翰逊）】}
所有Q输出的频率都是$f_{clk}/8$，但相位不同。

\textbf{【与二进制计数器对比】}
\begin{itemize}
  \item 二进制Q0：每2拍翻转→$f_{clk}/2$
  \item 二进制Q1：每4拍翻转→$f_{clk}/4$
  \item 二进制Q2：每8拍翻转→$f_{clk}/8$
  \item 约翰逊任意Q：每8拍翻转→$f_{clk}/8$（所有Q频率相同）
\end{itemize}

\end{solution}


% ============================================================
\section{题14：任意N位环形计数器模板}

\textbf{题目}：写出N位环形计数器的通用VHDL模板。

\begin{solution}[完整解答]

\textbf{【通用模板】}
\begin{lstlisting}
entity ring_counter is
  generic (N : integer := 4);
  port (clk, rst_n : in std_logic; q : out std_logic_vector(N-1 downto 0));
end ring_counter;

architecture behavioral of ring_counter is
  signal reg : std_logic_vector(N-1 downto 0);
begin
  process(clk, rst_n)
  begin
    if rst_n = '0' then
      reg <= (N-1 => '1', others => '0');  -- 只有MSB=1
    elsif rising_edge(clk) then
      reg <= reg(0) & reg(N-1 downto 1);  -- 右移+反馈
    end if;
  end process;
  q <= reg;
end behavioral;
\end{lstlisting}

\textbf{【参数化】}N可以是任意值。N=4→4个状态。N=8→8个状态。N个DFF→N个状态。

\end{solution}


% ============================================================
\section{题15：环形计数器类统一秒杀}

\begin{quickkill}[环形计数器秒杀总结]
\textbf{普通环形计数器：}
\begin{itemize}
  \item VHDL: \texttt{reg <= reg(0) \& reg(N-1 downto 1)}
  \item N个DFF → N个状态
  \item 初始：恰好一个1（One-hot）
  \item 每拍：1向右移一位（通过反馈形成循环）
\end{itemize}

\textbf{约翰逊（扭环形）计数器：}
\begin{itemize}
  \item VHDL: \texttt{reg <= (not reg(0)) \& reg(N-1 downto 1)}
  \item N个DFF → 2N个状态
  \item 初始：全0（0000）
  \item 规律：1从左填满，0从右清空
\end{itemize}

\textbf{口诀：}
\begin{itemize}
  \item 环形 = 1在环里跑，跑N步回原点
  \item 约翰逊 = 1填满再清空，跑2N步回原点
\end{itemize}
\end{quickkill}
